\chapter{Cache Memories}\label{ch:auto-004}


\section{Memory Technology and the Hierarchy}\label{sec:auto-019}
Modern cache design begins with a technology reality: fast memory is expensive and dense memory is slower.
\textit{Static random-access memory (SRAM)} cells use more transistors per bit, provide lower access latency, and do not require refresh, so SRAM is the natural choice for on-chip caches.
\textit{Dynamic random-access memory (DRAM)} cells are denser and cheaper per bit but slower and refresh-driven, so DRAM is the natural choice for main memory capacity.
This is why systems are built as hierarchies rather than as one uniform memory.

A useful mental model is that each level trades speed for size as distance from the core increases.
L1 and often L2 are SRAM-based and optimized for very low hit latency.
Lower levels and main memory prioritize capacity and cost.
Architecture work in Chapter 3 is mostly about exploiting this hierarchy so that accesses are absorbed by fast levels as often as possible.

At the circuit level, one practical limiter is wire delay.
Longer bit lines have higher effective capacitance and resistance, so sense amplifiers need more time to resolve values.
This is one reason very large single-level caches become slow: making arrays taller increases bit-line delay even if logic depth is unchanged.


\section{Typical Hierarchy Levels}\label{sec:auto-020}
In mainstream systems, the first level is a small L1 cache optimized for minimum hit time.
A larger L2 cache is usually next, trading some latency for much lower miss ratio.
Below L2 sits physical memory (DRAM), which is much larger but much slower.
Beneath DRAM is virtual-memory backing store (for example solid-state drive (SSD)), managed by the operating system through paging.

From the central processing unit (CPU) perspective, each lower level is accessed only when the upper level misses.
This staged access is exactly why average-access-time analysis composes recursively across levels.


\section{Locality: Why Caches Work}\label{sec:auto-021}
Cache behavior is dominated by \textit{locality}.
\textit{Temporal locality} means that if data is used now, it is likely to be used again soon.
\textit{Spatial locality} means that if one address is accessed, nearby addresses are likely to follow.
\textit{Sequential locality} is a common special case of spatial locality in which accesses proceed in order, such as instruction fetch or array traversal.

These properties justify block-based transfer.
Instead of moving one byte at a time from DRAM, the system moves an entire cache block, betting that adjacent bytes will be useful soon.
The same properties also justify replacement policies that retain recently used data and evict older candidates.


\section{Cache Geometry and Address Partitioning}\label{sec:auto-022}
To reason cleanly about cache structure, use the triplet $(C,B,S)$ (Eq.~\eqref{eq:auto-018}):
\begin{equation}\label{eq:auto-018}
    \text{cache capacity}=2^C\ \text{bytes},\qquad
    \text{block size}=2^B\ \text{bytes},\qquad
    \text{associativity}=2^S\ \text{ways}.
\end{equation}

\paragraph{Address fields in $(C,B,S)$ notation}\label{par:auto-012}
In this notation, $S$ is \textbf{not} an extra address field beyond tag/index/offset.
Instead, $S$ changes the number of sets and therefore the index width.
\begin{equation}\label{eq:auto-019}
    \text{sets}=\frac{2^C}{2^B\cdot 2^S}=2^{C-B-S}.
\end{equation}
Total line count follows directly from sets and ways (Eq.~\eqref{eq:auto-086}):
\begin{equation}\label{eq:auto-086}
    \text{lines}=\text{sets}\cdot\text{ways}=2^{C-B-S}\cdot 2^S=2^{C-B}.
\end{equation}
For direct-mapped caches ($S=0$), $\text{ways}=1$ so $\text{lines}=\text{sets}$.
For set-associative caches ($S>0$), $\text{lines}>\text{sets}$ because there are multiple ways per set.
\begin{equation}\label{eq:auto-020}
    |\text{Offset}|=B,\qquad
    |\text{Index}|=\log_2(\text{sets})=C-B-S.
\end{equation}
Because index plus offset consume $(C-S)$ low-order bits (Eq.~\eqref{eq:auto-021}):
\begin{equation}\label{eq:auto-021}
    |\text{Index}|+|\text{Offset}|=(C-B-S)+B=C-S.
\end{equation}
For an address width of $AW$ bits (Eq.~\eqref{eq:auto-022}):
\begin{equation}\label{eq:auto-022}
    |\text{Tag}|=AW-(|\text{Index}|+|\text{Offset}|)=AW-(C-S).
\end{equation}
So the split remains (Eq.~\eqref{eq:auto-023}):
\begin{equation}\label{eq:auto-023}
    \text{Address}=[\text{Tag}\mid\text{Index}\mid\text{Offset}],
\end{equation}
with
\begin{equation}\label{eq:auto-024}
    |\text{Offset}|=B,\qquad |\text{Index}|=C-B-S,\qquad |\text{Tag}|=AW-(C-S).
\end{equation}

\paragraph{Where the $S$ bits ``go''}\label{par:auto-013}
Higher associativity means fewer sets, so index width decreases by $S$ relative to a direct-mapped organization with the same $C$ and $B$.
That is why $S$ appears in index/tag equations but is not carried as an address field.

\paragraph{What selects the way?}\label{par:auto-014}
The way is selected by hardware after tag comparison, not by an explicit address field:
\begin{enumerate}
    \item Index selects one set.
    \item Tags for all $2^S$ ways in that set are compared in parallel.
    \item Match logic identifies the hitting way (if any).
    \item Internal way-select control drives the data mux (conceptually an $S$-bit code since $\log_2(2^S)=S$).
\end{enumerate}

\paragraph{Quick example: 32\,KB, 32\,B blocks, 4-way}\label{par:auto-015}
With $(C,B,S)=(15,5,2)$ (Eq.~\eqref{eq:auto-025}):
\begin{equation}\label{eq:auto-025}
    |\text{Index}|=C-B-S=15-5-2=8.
\end{equation}
For the same $C,B$ but direct-mapped ($S=0$) (Eq.~\eqref{eq:auto-026}):
\begin{equation}\label{eq:auto-026}
    |\text{Index}|=15-5-0=10.
\end{equation}
So moving to 4-way associativity reduces index width by 2 bits (equivalently reduces sets by $4\times$).

\paragraph{32-bit address specialization}\label{par:auto-016}
If $AW=32$ (Eq.~\eqref{eq:auto-027}):
\begin{equation}\label{eq:auto-027}
    |\text{Offset}|=B,\qquad
    |\text{Index}|=C-B-S,\qquad
    |\text{Tag}|=32-(C-S)=32-C+S.
\end{equation}

\paragraph{Worked example (metadata sizing)}\label{par:auto-017}
Consider a 64-bit machine with $(C,B,S)=(13,5,0)$, an 8\,KB direct-mapped cache with 32-byte blocks.
The cache has $2^{13-5}=256$ blocks and $2^{13-5-0}=256$ sets.
Tag width is $64-(13-0)=51$ bits.
If each block stores one valid bit and one dirty bit, metadata per block is 53 bits.
Total tag-store plus metadata is (Eq.~\eqref{eq:auto-028}):
\begin{equation}\label{eq:auto-028}
    256 \times 53 = 13{,}568\ \text{bits}.
\end{equation}
If $B$ is reduced to 3 while $C$ and $S$ are fixed, the number of blocks rises to $2^{10}=1024$, so tag overhead increases to $1024\times53=54{,}272$ bits.
This illustrates an important Chapter 3 point: smaller blocks can increase tag overhead significantly.

\paragraph{Concrete Tag-Store Accounting}\label{par:auto-017a}
The quickest reliable procedure is:
\begin{enumerate}
    \item determine the number of lines from $2^{C-B}$,
    \item determine tag width from the relevant physical or virtual address width,
    \item add only the metadata bits the problem explicitly asks for, and
    \item multiply per-line metadata by line count.
\end{enumerate}
For example, suppose a cache is \textit{write-back, write-allocate} and \textit{physically indexed, physically tagged} with $(C,B,S)=(18,7,3)$, 64-bit virtual addresses, 4\,KB pages, and 32\,GB of physical memory.
The physical address width is 35 bits because $32$\,GB $=2^{35}$ bytes.
The cache therefore has $2^{18-7}=2^{11}=2048$ lines, and the physical tag width is
\begin{equation}\label{eq:auto-028a}
    |\text{Tag}| = 35-(18-3)=20\ \text{bits}.
\end{equation}
Because the level is write-back, write-allocate, each line also needs one valid bit and one dirty bit, so per-line tag-store plus required metadata is 22 bits.
Total tag-store size is therefore
\begin{equation}\label{eq:auto-028b}
    2048 \times (20+1+1)=45{,}056\ \text{bits}.
\end{equation}
If the problem says to ignore replacement metadata such as LRU state, then it should not be added.

\paragraph{Address partitioning example}\label{par:auto-018}
Consider the same $(13,5,0)$ 64-bit cache.
For the virtual address \texttt{0000\_1234\_ABCD\_5678}:
\textbf{Offset (5 bits):} the lowest 5 bits of \texttt{78} (\texttt{0111\_1000}) are \texttt{11000}, which is \texttt{18} (decimal 24).
\textbf{Index (8 bits):} bits 5 through 12 come from hex digits \texttt{5}, \texttt{6}, and \texttt{7}.
In binary these digits are \texttt{0101}, \texttt{0110}, and \texttt{0111}; extracting bits 5--12 yields \texttt{1011\_0011} (binary), which is \texttt{B3}.
\textbf{Tag (51 bits):} the remaining upper bits are the tag.
This direct bit-mapping is the fundamental operation performed by hardware on every memory access.
It illustrates how physical cache structure determines exactly which address bits are used for indexing and comparison.


\section{Direct-Mapped, Set-Associative, and Fully Associative}\label{sec:auto-023}
A \textit{direct-mapped cache} has one way per set, so each block address maps to exactly one location.
Lookup is simple and fast, but conflict misses can be severe.
A \textit{set-associative cache} allows each index to hold multiple candidate blocks; it reduces conflicts at the cost of parallel tag comparisons and a way-select multiplexer.
A \textit{fully associative cache} removes index constraints entirely and minimizes conflict misses, but hardware search complexity and power become substantial for large capacities.

Under $(C,B,S)$ notation, direct-mapped means $S=0$.
Fully associative means one set, so $C-B-S=0$ and therefore $S=C-B$.
This distinction is critical, because confusing these two conditions leads to incorrect address math and wrong simulator behavior.

\paragraph{How each cache is constructed}\label{par:auto-019}
Direct-mapped construction uses one data line and one tag entry per index.
The index decoder selects one row, then one tag compare and one word select complete the access.
A $2^S$-way set-associative cache replicates tag and data entries $2^S$ times per index; all tags in the selected set are compared in parallel, and a way-select mux picks the winning line.
A fully associative cache has a single set containing all lines, so every access performs an associative search across the entire tag store and relies heavily on replacement metadata.


\section{Hit, Miss, and Average Access Time}\label{sec:auto-024}
For a cache hit, access latency is the \textit{hit time} $HT$.
For a miss, the processor first spends hit-time logic discovering the miss, then pays \textit{miss penalty} $MP$ to fetch from the next level.
The \textit{average access time} is (Eq.~\eqref{eq:auto-029}):
\begin{equation}\label{eq:auto-029}
    AAT = HT + MR \cdot MP,
\end{equation}
where $MR$ is \textit{miss ratio}.
In Chapter 3 analysis, Eq.~\eqref{eq:auto-029} is the cache analogue of the CPU-time equation in Eq.~\eqref{eq:auto-002}.

If multiple workloads share one machine and workload $i$ accounts for fraction $w_i$ of the accesses, then weighted average access time is
\begin{equation}\label{eq:auto-029a}
    AAT_{weighted}=\sum_i w_i \, AAT_i,
\end{equation}
with $\sum_i w_i=1$.
This weighted form is the correct way to compare cache organizations against a non-uniform workload mix.

Hit time is not a single gate delay.
It includes row decoding, word-line drive, bit-line sensing, tag comparison, way selection, and word/byte selection.
Cache-geometry choices alter these terms.
Taller arrays hurt bit-line delay, wider arrays increase mux and comparator complexity, and higher associativity increases tag-compare and selection cost.

\paragraph{Which change affects which AAT term?}\label{par:auto-018a}
When a problem asks which of $HT$, $MR$, or $MP$ changes, the most useful first-pass mapping is usually:
\begin{enumerate}
    \item \textbf{data-layout or tiling changes} primarily move \textit{miss ratio} because they change locality and conflict behavior,
    \item \textbf{VIPT versus PIPT} primarily moves \textit{hit time} because it changes whether TLB lookup can overlap cache indexing, and
    \item \textbf{write buffers} primarily reduce effective \textit{miss penalty} or visible lower-level service time on writes by decoupling cache progress from memory completion.
\end{enumerate}
This classification is approximate rather than absolute, but it is the right first-pass answer for most design analyses.


\section{The Three C's of Misses}\label{sec:auto-025}
Misses are classed as \textit{compulsory}, \textit{capacity}, or \textit{conflict}.
Compulsory misses are first-touch misses and are unavoidable unless data is brought in before first demand access.
Capacity misses happen because the working set exceeds what the cache can hold over the reuse interval.
Conflict misses happen when mapping restrictions force evictions even though total capacity would have been enough under freer placement.

In practice, this taxonomy provides a diagnosis strategy.
If misses are mainly conflict-driven, associativity or victim structures can help.
If misses are mainly capacity-driven, larger effective storage or better blocking can help.
If misses are mainly compulsory, prefetching and layout changes are primary tools.

\paragraph{Sizing for a repeating working set}\label{par:auto-019a}
For a repeating access pattern, the smallest cache that avoids capacity misses must hold every block that remains live across the reuse interval.
If $N_{live}$ distinct blocks must stay resident until reuse, then the cache needs at least $N_{live}$ lines.
That condition is necessary but not always sufficient in a direct-mapped or low-associativity cache, because several of those live blocks may still collide onto the same set.
This is why trace analysis must separate ``enough total lines'' from ``enough placement freedom.''


\section{Replacement Policies and Their Tradeoffs}\label{sec:auto-026}
When a set is full and a miss occurs, one resident block must be evicted.
\textit{Belady's optimal policy} is a theoretical lower bound that requires future knowledge and is not implementable directly.
\textit{Least recently used (LRU)} approximates this by evicting the least recently touched entry, and usually performs well under temporal locality.
\textit{First in, first out (FIFO)} ignores reuse and can evict hot lines.
\textit{Random} is simple and sometimes surprisingly competitive but unpredictable.
\textit{Least frequently used (LFU)} tracks frequency but can lag phase changes.

For high associativity, exact LRU state grows costly.
Practical designs use approximations, tree-based encodings, or hybrid insertion policies.
A notable policy distinction is insertion behavior: a line can be installed as \textit{most recently used (MRU)} (default in many demand-miss paths) or near LRU (useful for speculative prefetch lines to reduce pollution).
Related insertion-policy abbreviations are \textit{most recently used insertion policy (MIP)} and \textit{least recently used insertion policy (LIP)}.

\paragraph{Simulating LRU with timestamps}\label{par:auto-019b}
One exact way to simulate LRU is to attach a ``last used'' timestamp to each line in a set.
On every hit, update the touched line's timestamp to the current access time.
On a miss, fill an invalid line if one exists; otherwise evict the valid line with the smallest timestamp because that line was touched least recently.
This method is conceptually simple and is often the most direct way to hand-simulate LRU on short traces.

\paragraph{Worked cache-trace decode}\label{par:auto-019bb}
Consider a cache with $(C,B,S)=(6,4,1)$.
Then the cache has $2^{6-4}=4$ lines, arranged as $2^{6-4-1}=2$ sets with 2 ways per set.
So the offset is 4 bits, the index is 1 bit, and the remaining upper bits form the tag.
Now consider the access sequence \texttt{0x24}, \texttt{0x34}, \texttt{0x2C}, \texttt{0x44}.
Because the offset is 4 bits, the block numbers are \texttt{0x2}, \texttt{0x3}, \texttt{0x2}, and \texttt{0x4}.
Using the low-order bit of the block number as the set index, those accesses map to set 0, set 1, set 0, and set 0, with tags 1, 1, 1, and 2, respectively.
The first two accesses are compulsory misses, the third is a hit on the already resident block for \texttt{0x24}, and the fourth is another miss because it introduces a different tag into set 0.
This is the basic hand-simulation loop: decode block, split into set and tag, check the resident ways, then update replacement state.

\paragraph{Worked timestamp-LRU trace}\label{par:auto-019d}
Consider a 2-way set with two accesses already resident:
\begin{center}
\begin{tabular}{|c|c|}
    \hline
    Tag & Timestamp \\
    \hline
    A & 4 \\
    \hline
    B & 7 \\
    \hline
\end{tabular}
\end{center}
If the next access is a hit on A at time 8, then A's timestamp becomes 8 while B remains at 7.
If the following access is a miss at time 9 and no invalid line exists, then B is evicted because it has the smaller timestamp and is therefore the least recently used line.
This is the mechanical rule to remember: every hit refreshes exactly one timestamp, and every full-set miss evicts the smallest timestamp.

\paragraph{Exact LRU with a bit matrix}\label{par:auto-019c}
Another exact representation for an $n$-way set is an $n\times n$ recency matrix.
Entry $M[i][j]=1$ means way $i$ is more recent than way $j$.
When way $k$ is accessed, set row $k$ to 1, set column $k$ to 0, and keep the diagonal in its fixed convention.
After the update, the least-recently used way is the one whose off-diagonal row bits are all 0, because every other way is marked as newer than it.

\paragraph{Worked bit-matrix update}\label{par:auto-019e}
For a 4-way set, suppose the current access is to way 2.
Then row 2 is set to 1 against the other ways and column 2 is cleared:
\begin{center}
\begin{tabular}{|c|cccc|}
    \hline
    & 0 & 1 & 2 & 3 \\
    \hline
    0 & -- & ? & 0 & ? \\
    1 & ? & -- & 0 & ? \\
    2 & 1 & 1 & -- & 1 \\
    3 & ? & ? & 0 & -- \\
    \hline
\end{tabular}
\end{center}
The question marks keep their previous values because the relative recency order among the untouched ways does not change.
After several accesses, the LRU way is the row whose off-diagonal entries are all 0.
That row is the one that every other way claims to be newer than.


\section{Write Policies}\label{sec:auto-027}
Write behavior has two independent decisions.
On a write hit, either update only the cache and defer memory update (\textit{write-back}), or update both cache and memory immediately (\textit{write-through}).
On a write miss, either fetch the block then write it (\textit{write-allocate}), or write to lower memory without installing in cache (\textit{write-no-allocate}).

Write-back reduces memory traffic when repeated writes target the same line, but requires a dirty bit and careful eviction handling.
Because dirty data can live in cache before eviction, abrupt power loss can drop recently written values that have not yet reached lower levels.
Write-through simplifies coherence visibility and durability behavior because lower memory is updated on each write hit, and it does not require dirty metadata at that cache level, but it increases memory bandwidth pressure.
The common pairings are write-back plus write-allocate for reuse-heavy data paths, and write-through plus write-no-allocate for streaming paths with little read-after-write reuse.

\paragraph{Reading write traces}\label{par:auto-020a}
Write traces are easiest to interpret by separating three questions.
First, does the access hit or miss in the current cache state?
Second, does the policy allocate the block on a miss?
Third, if a valid victim is evicted, is it dirty?
In a write-back, write-allocate cache, a write miss usually installs the block and marks it dirty immediately.
Each later dirty eviction contributes one writeback event.
Dirty lines that remain in the cache at the end of the trace are not counted as writebacks unless the problem explicitly asks for a final flush.

\paragraph{Worked dirty-eviction trace}\label{par:auto-020b}
Consider a write-back, write-allocate cache whose relevant set can hold only two lines.
Suppose the set is initially empty and the access sequence is
\[
W(A),\ W(B),\ W(A),\ R(C),
\]
with $C$ mapping to the same full set as $A$ and $B$.
The first write to $A$ is a miss, so $A$ is allocated and marked dirty.
The first write to $B$ does the same for $B$.
The second write to $A$ is a hit and leaves $A$ dirty.
When $C$ is later accessed, one of the dirty lines $A$ or $B$ must be evicted; that eviction produces a writeback because the victim holds modified data that has not yet reached the next level.

\paragraph{Valid and dirty bits: meaning, need, and size impact}\label{par:auto-020}
Each cache line typically carries at least one \textit{valid bit}.
Valid$=1$ means the tag/data entry contains a meaningful memory block; valid$=0$ means the entry must not hit, even if tag bits happen to match after reset or replacement.
A \textit{dirty bit} is needed when that cache level uses \textit{write-back}: dirty$=1$ means the line was modified in cache and must be written to the next level before eviction; dirty$=0$ means eviction can discard the line without write-back traffic.
In a pure write-through level, dirty bits are not required because lower memory is updated on every write hit.

If line count is $N_{lines}=2^{C-B}$, valid-bit storage is $N_{lines}$ bits.
Dirty-bit storage is another $N_{lines}$ bits only when write-back is enabled at that level.
Relative to data-array size $2^C$ bytes, overhead from these two bits alone is (Eq.~\eqref{eq:auto-030}):
\begin{equation}\label{eq:auto-030}
    \text{overhead fraction}=\frac{n_v+n_d}{2^{B+3}},
\end{equation}
where $n_v=1$ (valid always) and $n_d\in\{0,1\}$ (dirty only for write-back).
Example: with 32-byte blocks ($B=5$), valid-only overhead is $1/256\approx0.39\%$, and valid+dirty overhead is $2/256\approx0.78\%$ (excluding tag, replacement, and coherence bits).


\section{Reducing Miss Ratio}\label{sec:auto-028}
The first lever is \textbf{cache size ($C$)}.
Increasing $C$ reduces capacity misses because more of the working set fits in the hierarchy.
However, larger caches increase hit time due to longer bit lines and decoder delays, and they consume more area and power.

The second lever is \textbf{associativity ($S$)}.
Increasing $S$ reduces conflict misses by providing more locations for a given index.
The benefit is most significant when moving from direct-mapped to 2-way or 4-way, after which returns diminish.
The cost includes increased hit time (due to the way-select multiplexer and parallel tag comparisons) and increased power consumption.

The third lever is \textbf{block size ($B$)}.
Increasing $B$ reduces compulsory misses by exploiting spatial locality (fetching more useful data per miss).
However, if $B$ becomes too large, it can increase conflict misses (by reducing the total number of blocks) and will always increase miss penalty (longer transfer time).
So the statement ``increasing block size cannot improve miss ratio'' is false.
With a fixed cache capacity, a larger block often improves miss ratio at first by capturing nearby useful data, but the improvement is not monotonic because fewer total lines remain available.

\paragraph{Hardware Prefetching}\label{par:auto-021}
Hardware prefetching targets compulsory and capacity effects by pulling data before demand.
A \textit{$+1$ prefetcher} fetches block $i+1$ when block $i$ is accessed.
A \textit{stride prefetcher} detects constant deltas (for example $i, i+k, i+2k$) and fetches the next predicted address.
A \textit{Markov prefetcher} uses a correlation table to learn non-linear transitions; if address $A$ is often followed by $B$ then $C$, the predictor learns and issues those future blocks when $A$ appears.
A \textit{hybrid prefetcher} combines a richer predictor such as Markov correlation with a simpler fallback such as next-block prefetching, so the machine can still make a reasonable guess when the correlation table has no useful entry.
To prevent \textit{cache pollution} (evicting useful data for speculative data that might not be used), many designs use a \textit{prefetch buffer}.
This is a small, separate storage area where prefetched lines sit until they are actually demanded by the CPU\@.
If demanded, they are moved into the main cache.

\paragraph{Prefetch insertion policy matters}\label{par:auto-021e}
Prefetched lines are often installed with a low-priority policy such as \textit{least recently used insertion policy (LIP)} rather than as most-recently used lines.
The goal is to avoid immediately evicting useful demand data for speculative lines that may never be touched.
This helps in set-associative caches, but it does not eliminate the placement limitations of a direct-mapped cache.
In a direct-mapped design, a $+1$ prefetch can still evict useful demand data if the prefetched block maps to the same line.

\paragraph{Trace-driven prefetch statistics}\label{par:auto-021c}
Detailed simulator outputs often distinguish three prefetch counters.
\textit{Prefetches issued} count speculative requests launched by the prefetcher.
A \textit{prefetch hit} means a later demand access found a line that was present only because of an earlier prefetch.
A \textit{prefetch miss}, sometimes called a useless prefetch, means the prefetched line was evicted before any demand access used it.
These counters matter because a prefetcher can lower demand misses and still be a poor design if it wastes bandwidth or pollutes the cache badly.


\section{Reducing Miss Penalty}\label{sec:auto-029}
Miss penalty is shaped strongly by transfer mechanics.
If a block has $2^B$ bytes and bus transfer width is $W$ bytes, then the number of sub-block transfers is $N_{sb}=2^B/W$.
With first-transfer startup latency $T_1$ and per-following transfer time $T_f$, block transfer time is (Eq.~\eqref{eq:auto-031}):
\begin{equation}\label{eq:auto-031}
    MP_{full} = T_1 + (N_{sb}-1)T_f.
\end{equation}
Early restart returns data to the CPU as soon as the requested sub-block arrives rather than after full-block completion.
If critical-word position is uniformly distributed and transfer order is fixed from sub-block 0 upward, average CPU-visible penalty becomes (Eq.~\eqref{eq:auto-032}):
\begin{equation}\label{eq:auto-032}
    MP_{ER} \approx T_1 + \frac{N_{sb}-1}{2}T_f.
\end{equation}
Critical-word-first improves further by requesting the demanded sub-block first, then wrapping to complete the rest in the background.
That benefit is only CPU-visible if the machine can resume execution as soon as the critical word arrives.
Without early restart or an equivalent partial-consume mechanism, critical-word-first alone does not reduce the stall seen by the processor, because the core would still wait for the full block to complete before continuing.

\paragraph{Non-uniform critical-word distributions}\label{par:auto-021a}
The uniform approximation in Eq.~\eqref{eq:auto-032} is only one special case.
If sub-block $k$ is demanded with probability $p_k$ and fixed-order transfer would deliver that word after $k$ follow-up transfers, then the early-restart penalty is
\begin{equation}\label{eq:auto-032a}
    MP_{ER}=\sum_{k=0}^{N_{sb}-1} p_k \left(T_1 + kT_f\right).
\end{equation}
Under the same simple model, critical-word-first with early restart reduces the CPU-visible penalty to approximately $T_1$ because the demanded sub-block is requested first and the rest of the block finishes in the background.

\paragraph{Worked miss-penalty comparison}\label{par:auto-021aa}
Consider a 32-byte block on an 8-byte bus, so the miss repairs 4 words.
Suppose each requested word transfer takes 70 ns and transfers are serialized.
Baseline repair waits for all 4 words, so the visible miss penalty is
\begin{equation}\label{eq:auto-032b}
    MP_{full}=4\times 70 = 280\ \text{ns}.
\end{equation}
With early restart under a uniform critical-word position, the demanded word arrives after 1, 2, 3, or 4 transfers with equal probability, so the average visible penalty is
\begin{equation}\label{eq:auto-032c}
    MP_{ER}=\frac{70+140+210+280}{4}=175\ \text{ns}.
\end{equation}
With critical-word-first plus early restart, the demanded word is requested first, so the processor-visible penalty returns to approximately one transfer time, 70 ns, while the rest of the block completes in the background.

Another major lever is adding lower cache levels.
If L1 misses probe L2 before DRAM, L1 penalty becomes L2 average access time (Eq.~\eqref{eq:auto-033}):
\begin{equation}\label{eq:auto-033}
    AAT_{L1} = HT_{L1} + MR_{L1}\cdot AAT_{L2}.
\end{equation}
This recursive form extends naturally to deeper hierarchies.

\paragraph{Worked two-level AAT example}\label{par:auto-021ac}
Suppose $HT_{L1}=1$ cycle, $MR_{L1}=0.08$, $HT_{L2}=8$ cycles, $MR_{L2,read}=0.25$, and $DRAM\_TIME=60$ cycles.
Then
\begin{equation}\label{eq:auto-033a}
    AAT_{L2}=8+0.25\times 60=23\ \text{cycles}.
\end{equation}
That lower-level average time becomes the miss penalty seen by L1, so
\begin{equation}\label{eq:auto-033b}
    AAT_{L1}=1+0.08\times 23=2.84\ \text{cycles}.
\end{equation}
The main reasoning pattern is recursive rather than additive-by-guess: first solve the lower level, then substitute that average behavior upward.

\paragraph{Inclusive policy: subset, not equality}\label{par:auto-021f}
If L2 is \textit{inclusive} of L1, then every line present in L1 must also be present in L2.
The reverse is not required: L2 can contain many lines that are not currently in L1.
So inclusivity means ``L1 is a subset of L2,'' not ``L1 and L2 contain exactly the same data.''

\paragraph{Why many L2 ratios are read-only}\label{par:auto-021d}
In a hierarchy where L1 uses write-back, write-allocate and L2 uses write-through, write-no-allocate, simulators often report separate L2 read and write traffic.
Read hits and read misses form a meaningful reuse ratio because reads probe the resident L2 contents.
Writes, by contrast, may simply reflect L1 writebacks or pass-through stores and may not allocate lines on a miss.
For that reason, many reports define L2 hit ratio and miss ratio over \textit{reads only} rather than over total L2 traffic.

Victim caches are also miss-penalty tools.
A small fully associative buffer captures lines evicted from the primary cache.
On a later conflict-driven re-reference, data can be swapped back quickly without going to DRAM\@.
Conceptually, a victim cache provides a small amount of flexible associativity exactly where pressure appears.

\paragraph{Victim cache versus L2}\label{par:auto-021b}
A victim cache and an L2 do not solve exactly the same problem.
A victim cache is usually the better value when misses are dominated by a small number of conflict patterns near the L1 boundary, because a few fully associative entries can undo repeated index collisions cheaply.
A larger L2 is usually the better value when misses reflect broader capacity pressure or many different miss streams, because it filters a much wider set of misses rather than only recently evicted L1 lines.
The right choice depends on whether the workload is mainly suffering from local conflicts or from a deeper shortage of effective capacity.

Write buffers reduce stalls by decoupling cache writes from lower-memory completion.
They are especially useful when a following miss would otherwise wait behind pending writes.
Correct implementation requires forwarding from the newest matching buffered write and careful ordering semantics.


\section{Reducing Hit Time}\label{sec:auto-030}
Hit-time optimization frequently conflicts with miss-rate optimization.
Smaller arrays and lower associativity usually improve hit latency.
However, overly aggressive simplification can raise miss ratio enough to hurt total performance.
In general, small caches and direct-mapped organizations are fast on the hit path, but they are often unacceptable unless miss ratio remains controlled by workload behavior or lower-level caches.

Address translation complicates this tradeoff.
A naive \textit{physically indexed, physically tagged} path can serialize translation lookaside buffer (TLB) lookup before cache access, increasing L1 hit latency.
\textit{Virtually indexed, physically tagged (VIPT)} caches overlap TLB translation with set indexing and preserve physical-tag correctness.

\paragraph{PIPT vs. VIPT timing view}\label{par:auto-074}
With virtual memory, the CPU generates a virtual address split into virtual page number (VPN) and page offset.
In a \textit{physically indexed, physically tagged (PIPT)} L1, both index and tag use physical address bits.
That means the TLB must produce the physical address before cache indexing can begin.
So a first-order timing model is (Eq.~\eqref{eq:auto-088}):
\begin{equation}\label{eq:auto-088}
    HT_{L1,\mathrm{PIPT}} \approx T_{TLB}+T_{L1,\mathrm{index+array+tag}}.
\end{equation}
In a \textit{VIPT} L1, indexing starts immediately using low-order virtual bits while the TLB is translating in parallel.
When the physical address arrives, the physical tag compare finalizes hit/miss.
A first-order model is (Eq.~\eqref{eq:auto-089}):
\begin{equation}\label{eq:auto-089}
    HT_{L1,\mathrm{VIPT}} \approx \max\!\left(T_{TLB},\,T_{L1,\mathrm{index+array}}\right)+T_{\mathrm{phys\_tag\_compare}}.
\end{equation}
The practical effect is that VIPT can hide much of TLB latency on the L1 hit path, while PIPT often exposes it directly.

\paragraph{Why the VIPT constraint exists}\label{par:auto-075}
For VIPT to work cleanly, index-plus-offset bits must fit within page-offset bits (Eq.~\eqref{eq:auto-034}):
\begin{equation}\label{eq:auto-034}
    C-S \le P,
\end{equation}
where $P$ is page-offset width in bits.
If $C-S>P$, then some index bits come from VPN bits.
Different virtual addresses that map to the same physical page (\textit{synonyms}) can then index different sets, allowing duplicated physical data in L1.
That creates alias-management problems: writes through one synonym may not immediately update the other cached copy, risking stale reads until extra coherence/synchronization machinery repairs it.
This constraint often forces minimum associativity for a given L1 size and page size.
For example, with $P=12$ (4\,KB pages) and $C=15$ (32\,KB L1), one needs $S\ge 3$, i.e., at least 8-way associativity.
If $S=2$ (4-way), then $C-S=13>12$, so VIPT indexing would extend beyond page offset and synonym risk appears.

\paragraph{Concrete VIPT Associativity Check}\label{par:auto-075a}
Suppose $C=18$ and page size is 4\,KB, so $P=12$.
To make the cache VIPT without indexing beyond the page offset, require
\begin{equation}\label{eq:auto-089a}
    C-S \le P \Rightarrow 18-S \le 12 \Rightarrow S \ge 6.
\end{equation}
So the minimum acceptable associativity field is $S=6$, which corresponds to $2^6=64$ ways.
The arithmetic is simple; the main thing to remember is that the constraint is on \emph{index plus offset}, which in $(C,B,S)$ notation is exactly $C-S$.

Write buffers can also improve effective hit behavior for writes by decoupling write completion from the immediate critical path and by forwarding recent write data when needed.
This prevents some write-after-read and read-after-write timing stalls from stretching apparent hit latency.


\section{Virtual Memory Technology, Terminology, and Analogy to Caching}\label{sec:auto-031}
Virtual-memory terminology mirrors cache terminology closely.
A virtual address is typically split into a \textit{virtual page number (VPN)} and page offset.
Translation maps VPN to a \textit{physical frame number (PFN)} through a page table.
A \textit{translation lookaside buffer (TLB)} is a cache of recent translations.
A missing translation in the TLB causes a page-table walk; if the page is not in DRAM, a page fault triggers operating system (OS) handling and secondary-storage transfer.

The analogy to caching is direct.
VPN behaves like a tag, page offset behaves like block offset, and the TLB behaves like a small associative cache of mapping metadata.
Likewise, DRAM behaves as a cache for backing store, while on-chip caches behave as caches for DRAM\@.
This layered view is useful because it applies one conceptual framework to both hardware cache hierarchy and OS-managed virtual memory.


\section{Compiler and Data-Layout Techniques}\label{sec:auto-032}
Software organization strongly affects cache performance.
\textbf{Compiler layout of instructions} places frequently executed basic blocks close together in memory to improve I-cache hit rate and reduce fetch stalls.
\textbf{Compiler-directed prefetching} issues non-binding prefetch instructions ahead of use, which is effective for predictable array accesses.
\textbf{Data-layout choices}, including structure-of-arrays (SoA), field reordering, and padding, can reduce cache-line waste and false sharing.

\textbf{Tiling (blocking)} re-organizes loops to process data in smaller cache-resident regions.
For example:
\begin{verbatim}
// Before tiling:
for (i=0; i<N; i++) for (j=0; j<N; j++) A[i][j] = B[j][i];
// After tiling (where b is the tile size):
for (ii=0; ii<N; ii+=b) for (jj=0; jj<N; jj+=b)
    for (i=ii; i<ii+b; i++) for (j=jj; j<jj+b; j++) A[i][j] = B[j][i];
\end{verbatim}

Choosing the tile factor is a working-set problem.
The largest profitable tile is usually the largest one whose active data footprint still fits in the effective cache space while avoiding severe mapping conflicts.
For direct-mapped or low-associativity caches, this means a tile can fail even when the byte count appears to fit, because the active blocks may still alias onto the same few sets.
The practical rule is therefore: first check footprint size, then check placement conflicts.
This is why tiled matrix multiplication often shows much better cache behavior than a naive matrix multiply, even when both perform the same arithmetic work.

\paragraph{Concrete Tile-Factor Estimate}\label{par:auto-075b}
Consider a cache with $(C,B,S)=(10,6,2)$, so total capacity is $2^{10}=1024$ bytes, block size is 64 bytes, and associativity is 4-way.
If each matrix element is 8 bytes and the tiled working set keeps $TF$ rows of 32 elements active, then one tiled row occupies
\begin{equation}\label{eq:auto-089b}
    32 \times 8 = 256\ \text{bytes}.
\end{equation}
A footprint estimate for the tile is therefore $256 \times TF$ bytes.
The largest tile that still fits in 1024 bytes is
\begin{equation}\label{eq:auto-089c}
    256 \times TF \le 1024 \Rightarrow TF \le 4.
\end{equation}
Here $TF=4$ is also structurally reasonable because each row occupies four 64-byte blocks, and the 4-way organization provides enough placement freedom to avoid immediate within-tile self-conflict.

\textbf{Loop fusion} combines independent loops that touch related data to improve temporal locality.
For example:
\begin{verbatim}
for (i=0; i<N; i++) a[i] = b[i] + 1; // loop 1
for (i=0; i<N; i++) c[i] = a[i] * 2; // loop 2
// Fused:
for (i=0; i<N; i++) { a[i] = b[i] + 1; c[i] = a[i] * 2; }
\end{verbatim}

In practice, these compiler and data-layout techniques frequently determine whether hardware cache features help or saturate.


\section{Chapter 3 Summary}\label{sec:auto-033}
Chapter 3 develops cache memory from first principles to practical optimization.
The chapter links SRAM/DRAM technology limits to hierarchical design, formalizes cache geometry with $(C,B,S)$ math, and uses AAT to connect structure to performance.
Misses are interpreted through the three C's, while replacement and write policies define behavioral tradeoffs.
Improvement strategies are organized by target term: miss ratio, miss penalty, or hit time.
Prefetching, victim caches, write buffers, multi-level hierarchies, simulator-style hierarchy and prefetch counters, and VIPT constraints show how modern designs balance speed, capacity, and correctness.
Finally, compiler and data-layout techniques complete the picture by showing that cache performance is a hardware-software co-design problem.


\section{Practice Questions}\label{sec:ch3-practice}
The following questions reinforce the main ideas of Chapter 3 and are intended for self-checking after the summary.
\begin{enumerate}
    \item A physically indexed, physically tagged cache has total capacity 32 KiB, block size 64 B, associativity 4-way, and 48-bit addresses. Compute the number of sets, the number of offset bits, the number of index bits, and the number of tag bits. Then compute the total tag-store size if each line stores a valid bit and a dirty bit in addition to the tag.
    \item For each of the following changes, state whether it primarily targets \emph{hit time}, \emph{miss ratio}, or \emph{miss penalty}: adding a victim cache, adding a write buffer, increasing associativity, increasing block size, and using critical-word-first with early restart.
    \item Explain the difference between critical-word-first and early restart. Why does critical-word-first by itself not guarantee a reduction in processor-visible miss penalty?
    \item A system uses 32-byte blocks, an 8-byte memory bus, and serialized memory transfers. A miss first incurs a 70 ns round trip for the requested word, and each additional word also requires 70 ns. The cache hit time is 10 ns and the hit rate is 15\%. Compute the average access time under three cases: baseline repair with no early use of returned data, early restart with a uniform probability of needing any of the four words in the block, and critical-word-first with early restart.
    \item Consider a $(C,B,S)=(6,4,2)$ cache with LRU replacement and the access stream \texttt{0x589e}, \texttt{0xa7a1}, \texttt{0x3767}, \texttt{0xcce2}, \texttt{0xa7aa}, \texttt{0xff4d}, \texttt{0x376c}, \texttt{0x5892}. For each access, determine the set, tag, hit or miss outcome, and final contents of the tag store after the full trace.
    \item A 4-way set uses an exact LRU bit-matrix implementation. Starting from a state in which way 0 is least recently used and way 3 is most recently used, update the matrix for the hit sequence $2\rightarrow 0\rightarrow 3\rightarrow 1\rightarrow 2$. Which way is least recently used at the end?
    \item A VIPT L1 cache has capacity 64 KiB, block size 64 B, and page size 4 KiB. Compute the minimum associativity required so that the set index fits entirely within the page offset.
    \item A miss stream repeatedly alternates between two blocks that map to the same direct-mapped L1 set. Would a small victim cache or a simple $+1$ prefetcher be the more natural first fix? Explain why.
    \item A tiled matrix kernel keeps three $b\times b$ tiles of double-precision data active at once. If the effective L1 capacity available to the tiles is 24 KiB, what is the largest integer tile size $b$ that satisfies the footprint constraint $3b^2\cdot 8 \le 24\text{ KiB}$?
    \item A two-level hierarchy has $HT_{L1}=1$ cycle, $MR_{L1}=0.08$, $HT_{L2}=8$ cycles, $MR_{L2,read}=0.25$, and $DRAM\_TIME=60$ cycles. Compute $AAT_{L2}$ and then $AAT_{L1}$.
    \item A write-back, write-allocate cache receives the sequence: write block A, write block B, write block A again, then access block C, which maps to the same full set as A and B and forces an eviction. Describe when dirty state is created and whether the eviction causes a writeback.
\end{enumerate}
